\documentclass[10pt,a4paper]{article}
\usepackage[utf8]{inputenc}
\usepackage[T1]{fontenc}
\usepackage[french]{babel}
\usepackage[margin=1.5cm]{geometry}
\usepackage{lmodern}
\usepackage{amsmath}
\usepackage{graphicx}
\usepackage{xcolor}
\usepackage{listings}
\usepackage{tcolorbox}
\tcbuselibrary{breakable} 
\usepackage{textcomp}
\usepackage{tikz}
\usetikzlibrary{shapes.geometric, arrows, positioning}

% Styles TikZ
\tikzstyle{startstop} = [rectangle, rounded corners, minimum width=3cm, minimum height=1cm,text centered, draw=black, fill=red!30, align=center]
\tikzstyle{io} = [trapezium, trapezium left angle=70, trapezium right angle=110, minimum width=3cm, minimum height=1cm, text centered, draw=black, fill=blue!30, align=center]
\tikzstyle{process} = [rectangle, minimum width=3cm, minimum height=1cm, text centered, draw=black, fill=orange!30, align=center]
\tikzstyle{decision} = [diamond, minimum width=3cm, minimum height=1cm, text centered, draw=black, fill=green!30, aspect=2, align=center]
\tikzstyle{arrow} = [thick,->,>=stealth]

% Couleurs
\definecolor{mainBlue}{HTML}{0D47A1}
\definecolor{softTeal}{HTML}{4DB6AC}
\definecolor{alertRed}{HTML}{D32F2F}
\definecolor{bgGray}{HTML}{F5F5F5}
\definecolor{codebg}{gray}{0.95}

\renewcommand{\familydefault}{\sfdefault}

% Configuration de la boîte d'exercice
\newtcolorbox{exoBox}[1]{
  colback=bgGray!50,
  colframe=mainBlue,
  fonttitle=\bfseries,
  coltitle=white,
  title=#1,
  rounded corners,
  arc=3mm,
  boxrule=1pt,
  left=5mm,
  right=5mm,
  top=3mm,
  bottom=3mm,
  breakable
}

% Style du code Python
\lstdefinestyle{pythonStyle}{
    backgroundcolor=\color{codebg},
    commentstyle=\color{green!60!black},
    keywordstyle=\color{blue},
    stringstyle=\color{purple},
    basicstyle=\ttfamily\small,
    breakatwhitespace=false,
    breaklines=true,
    captionpos=b,
    keepspaces=true,
    numbers=left,
    numbersep=5pt,
    showspaces=false,
    showstringspaces=false,
    showtabs=false,
    tabsize=4,
    language=Python,
    literate={é}{{\'e}}1 {à}{{\`a}}1 {è}{{\`e}}1 {ê}{{\^e}}1 {ç}{{\c c}}1 {ù}{{\`u}}1 {°}{{\textdegree}}1
}
\lstset{style=pythonStyle}

\pagestyle{empty}

\begin{document}

\begin{center}
    \huge\textbf{Initiation à la programmation} \\
    \vspace{0.2cm}
    \large Travail en autonomie \\
    \vspace{0.5cm}
    \large Maël Mignard \\
    \large 26 janvier 2026
\end{center}
\vspace{1cm}

\begin{exoBox}{Exercice 1 : Calcul du discriminant}
\subsection*{Organigramme}
\begin{center}
\resizebox{\textwidth}{!}{
    \begin{tikzpicture}[node distance=2cm]
    \node (start) [startstop] {Début};
    \node (in1) [io, below of=start] {Entrée: a, b, c};
    \node (pro1) [process, below of=in1] {delta = $b^2 - 4ac$};
    \node (dec1) [decision, below of=pro1, yshift=-1cm] {delta > 0 ?};
    \node (pro2a) [process, below of=dec1, yshift=-2cm, xshift=-4cm] {$x_1 = (-b - \sqrt{\text{delta}}) / (2a)$\\$x_2 = (-b + \sqrt{\text{delta}}) / (2a)$};
    \node (out1) [io, below of=pro2a] {Sortie: $x_1, x_2$};
    \node (dec2) [decision, right of=dec1, xshift=4cm] {delta == 0 ?};
    \node (pro2b) [process, below of=dec2, yshift=-2cm] {$x = -b / (2a)$};
    \node (out2) [io, below of=pro2b] {Sortie: x};
    \node (out3) [io, right of=dec2, xshift=4cm] {Sortie: "Solutions complexes"};
    \node (stop) [startstop, below of=out1, yshift=-1cm] {Fin};

    \draw [arrow] (start) -- (in1);
    \draw [arrow] (in1) -- (pro1);
    \draw [arrow] (pro1) -- (dec1);
    \draw [arrow] (dec1) -- node[anchor=east] {oui} (pro2a);
    \draw [arrow] (dec1) -- node[anchor=south] {non} (dec2);
    \draw [arrow] (pro2a) -- (out1);
    \draw [arrow] (out1) -- (stop);
    \draw [arrow] (dec2) -- node[anchor=east] {oui} (pro2b);
    \draw [arrow] (dec2) -- node[anchor=south] {non} (out3);
    \draw [arrow] (pro2b) -- (out2);
    \draw [arrow] (out2) |- (stop);
    \draw [arrow] (out3) |- (stop);
    \end{tikzpicture}
}
\end{center}

\subsection*{Programme Python}
\noindent
\begin{minipage}{\linewidth}
\begin{lstlisting}
import math

print("Calcul des racines d'une équation du second degré")
a = float(input("Entrez le coefficient a : "))
b = float(input("Entrez le coefficient b : "))
c = float(input("Entrez le coefficient c : "))

delta = b**2 - 4*a*c

if delta > 0:
    print(f"Le discriminant est positif (delta = {delta}).")
    x1 = (-b - math.sqrt(delta)) / (2*a)
    x2 = (-b + math.sqrt(delta)) / (2*a)
    print(f"Les deux racines réelles sont : x1 = {x1} et x2 = {x2}")
elif delta == 0:
    print(f"Le discriminant est nul (delta = {delta}).")
    x = -b / (2*a)
    print(f"La racine double est : x = {x}")
else:
    print(f"Le discriminant est négatif (delta = {delta}).")
    # Partie réelle et imaginaire pour l'affichage
    real_part = -b / (2*a)
    imag_part = math.sqrt(-delta) / (2*a)
    print("Les deux racines complexes sont :")
    print(f"x1 = {real_part} - {imag_part}i")
    print(f"x2 = {real_part} + {imag_part}i")
\end{lstlisting}
\end{minipage}

\subsection*{Exemples d\'exécution}
\begin{verbatim}
# Cas delta > 0 (a=1, b=-3, c=2)
...
Le discriminant est positif (delta = 1.0).
Les deux racines réelles distinctes sont : x1 = 1.0 et x2 = 2.0
\end{verbatim}
\end{exoBox}
\vspace{0.5cm}

\begin{exoBox}{Exercice 2 : Conversion de température}
\subsection*{Organigramme}
\begin{center}
\resizebox{\textwidth}{!}{
    \begin{tikzpicture}[node distance=2cm]
    \node (start) [startstop] {Début};
    \node (in1) [io, below of=start] {Entrée: temp\_c};
    \node (in2) [io, below of=in1] {Entrée: choix (1 ou 2)};
    \node (dec1) [decision, below of=in2, yshift=-1cm] {choix == 1 ?};
    \node (pro1) [process, below of=dec1, yshift=-2cm, xshift=-3cm] {temp\_f = temp\_c * 9/5 + 32};
    \node (out1) [io, below of=pro1] {Sortie: temp\_f};
    \node (dec2) [decision, right of=dec1, xshift=4cm] {choix == 2 ?};
    \node (pro2) [process, below of=dec2, yshift=-2cm] {temp\_k = temp\_c + 273.15};
    \node (out2) [io, below of=pro2] {Sortie: temp\_k};
    \node (out3) [io, right of=dec2, xshift=4cm] {Sortie: "Choix invalide"};
    \node (stop) [startstop, below of=out1, yshift=-1cm] {Fin};

    \draw [arrow] (start) -- (in1);
    \draw [arrow] (in1) -- (in2);
    \draw [arrow] (in2) -- (dec1);
    \draw [arrow] (dec1) -- node[anchor=east] {oui} (pro1);
    \draw [arrow] (dec1) -- node[anchor=south] {non} (dec2);
    \draw [arrow] (pro1) -- (out1);
    \draw [arrow] (out1) -- (stop);
    \draw [arrow] (dec2) -- node[anchor=east] {oui} (pro2);
    \draw [arrow] (dec2) -- node[anchor=south] {non} (out3);
    \draw [arrow] (pro2) -- (out2);
    \draw [arrow] (out2) |- (stop);
    \draw [arrow] (out3) |- (stop);
    \end{tikzpicture}
}
\end{center}
\subsection*{Programme Python}
\noindent
\begin{minipage}{\linewidth}
\begin{lstlisting}
print("Conversion de température")
temp_c = float(input("Entrez la température en degrés Celsius : "))
print("Choisissez le mode de conversion :")
print("1. Convertir en Fahrenheit")
print("2. Convertir en Kelvin")
choix = int(input("Votre choix (1 ou 2) : "))

if choix == 1:
    temp_f = temp_c * 9/5 + 32
    print(f"{temp_c}°C équivaut à {temp_f}°F.")
elif choix == 2:
    temp_k = temp_c + 273.15
    print(f"{temp_c}°C équivaut à {temp_k}K.")
else:
    print("Erreur : Choix invalide. Veuillez entrer 1 ou 2.")
\end{lstlisting}
\end{minipage}

\subsection*{Exemples d\'exécution}
\begin{verbatim}
# Conversion en Fahrenheit (C=25, choix=1)
Conversion de température
Entrez la température en degrés Celsius : 25
Choisissez le mode de conversion :
1. Convertir en Fahrenheit
2. Convertir en Kelvin
Votre choix (1 ou 2) : 1
25.0°C équivaut à 77.00°F.

# Conversion en Kelvin (C=25, choix=2)
Conversion de température
Entrez la température en degrés Celsius : 25
Choisissez le mode de conversion :
1. Convertir en Fahrenheit
2. Convertir en Kelvin
Votre choix (1 ou 2) : 2
25.0°C équivaut à 298.15K.

# Choix invalide (C=25, choix=3)
Conversion de température
Entrez la température en degrés Celsius : 25
Choisissez le mode de conversion :
1. Convertir en Fahrenheit
2. Convertir en Kelvin
Votre choix (1 ou 2) : 3
Erreur : Choix invalide. Veuillez entrer 1 ou 2.

\end{verbatim}
\end{exoBox}
\vspace{0.5cm}

\begin{exoBox}{Exercice 3 : Calcul des aires de figures géométriques}
\subsection*{Organigramme}
\begin{center}
\resizebox{\textwidth}{!}{
    \begin{tikzpicture}[node distance=2cm]
    \node (start) [startstop] {Début};
    \node (in1) [io, below of=start] {Entrée: choix (1-4)};
    \node (dec1) [decision, below of=in1, yshift=-1cm] {choix?};

    \node (pro1) [process, below of=dec1, yshift=-3.5cm, xshift=-6cm] {Entrée: longueur, largeur\\aire = longueur $\times$ largeur};
    \node (out1) [io, below of=pro1] {Sortie: aire};

    \node (pro2) [process, below of=dec1, yshift=-3.5cm, xshift=-2cm] {Entrée: rayon\\aire = $\pi \times \text{rayon}^2$};
    \node (out2) [io, below of=pro2] {Sortie: aire};

    \node (pro3) [process, below of=dec1, yshift=-3.5cm, xshift=2cm] {Entrée: base, hauteur\\aire = 0.5 $\times$ base $\times$ hauteur};
    \node (out3) [io, below of=pro3] {Sortie: aire};

    \node (pro4) [process, below of=dec1, yshift=-3.5cm, xshift=6cm] {Entrée: b1,b2,h\\aire=0.5 $\times$ (b1+b2) $\times$ h};
    \node (out4) [io, below of=pro4] {Sortie: aire};

    \node (stop) [startstop, below of=out2, yshift=-2cm] {Fin};

    \draw [arrow] (start) -- (in1);
    \draw [arrow] (in1) -- (dec1);

    \draw [arrow] (dec1) -- node[anchor=south] {1} (pro1);
    \draw [arrow] (dec1) -- node[anchor=south] {2} (pro2);
    \draw [arrow] (dec1) -- node[anchor=south] {3} (pro3);
    \draw [arrow] (dec1) -- node[anchor=south] {4} (pro4);

    \draw [arrow] (pro1) -- (out1);
    \draw [arrow] (pro2) -- (out2);
    \draw [arrow] (pro3) -- (out3);
    \draw [arrow] (pro4) -- (out4);

    \draw [arrow] (out1) -- (stop);
    \draw [arrow] (out2) -- (stop);
    \draw [arrow] (out3) -- (stop);
    \draw [arrow] (out4) -- (stop);
    \end{tikzpicture}
}
\end{center}
\subsection*{Programme Python}
\noindent
\begin{minipage}{\linewidth}
\begin{lstlisting}
import math

print("Calcul d'aires de figures géométriques")
print("1. Rectangle | 2. Cercle | 3. Triangle | 4. Trapèze")
choix = int(input("Votre choix (1-4) : "))

if choix == 1:
    L = float(input("Entrez la longueur : "))
    l = float(input("Entrez la largeur : "))
    print(f"L'aire du rectangle est : {L * l}")
elif choix == 2:
    r = float(input("Entrez le rayon : "))
    print(f"L'aire du cercle est : {round(math.pi * r**2, 2)}")
elif choix == 3:
    b = float(input("Entrez la base : "))
    h = float(input("Entrez la hauteur : "))
    print(f"L'aire du triangle est : {0.5 * b * h}")
elif choix == 4:
    b1 = float(input("Entrez la petite base : "))
    b2 = float(input("Entrez la grande base : "))
    h = float(input("Entrez la hauteur : "))
    print(f"L'aire du trapèze est : {0.5 * (b1 + b2) * h}")
else:
    print("Choix invalide.")
\end{lstlisting}
\end{minipage}

\subsection*{Exemples d\'exécution}
\begin{verbatim}
# Cas 1: Rectangle (longueur=10, largeur=5)
...
L'aire du rectangle est : 50.0
\end{verbatim}
\end{exoBox}
\vspace{0.5cm}

\begin{exoBox}{Exercice 4 : Calculatrice d'opérations de base}
\subsection*{Organigramme}
\begin{center}
\resizebox{\textwidth}{!}{
    \begin{tikzpicture}[node distance=2cm]
    \node (start) [startstop] {Début};
    \node (in1) [io, below of=start] {Entrée: op, x, y};
    \node (dec1) [decision, below of=in1, yshift=-1cm] {op?};

    \node (pro_add) [process, below of=dec1, yshift=-3cm, xshift=-9cm] {res = x + y};
    \node (pro_sub) [process, below of=dec1, yshift=-3cm, xshift=-5cm] {res = x - y};
    \node (pro_mul) [process, below of=dec1, yshift=-3cm, xshift=-1cm] {res = x * y};
    \node (pro_pow) [process, below of=dec1, yshift=-3cm, xshift=3cm] {res = x ** y};

    \node (dec_div) [decision, below of=dec1, yshift=-3cm, xshift=7cm] {y == 0?};
    \node (pro_div) [process, below of=dec_div, yshift=-2cm] {res = x / y};
    \node (out_div_err) [io, right of=dec_div, xshift=3cm] {Sortie: "Div par zéro"};

    \node (out) [io, below of=pro_sub, yshift=-2cm, xshift=2cm] {Sortie: res};
    \node (stop) [startstop, below of=out] {Fin};

    \draw [arrow] (start) -- (in1);
    \draw [arrow] (in1) -- (dec1);

    \draw [arrow] (dec1) -- node[anchor=south] {+} (pro_add);
    \draw [arrow] (dec1) -- node[anchor=south] {-} (pro_sub);
    \draw [arrow] (dec1) -- node[anchor=south] {*} (pro_mul);
    \draw [arrow] (dec1) -- node[anchor=south] {**} (pro_pow);
    \draw [arrow] (dec1) -- node[anchor=south] {/} (dec_div);

    \draw [arrow] (pro_add) |- (out);
    \draw [arrow] (pro_sub) -- (out);
    \draw [arrow] (pro_mul) |- (out);
    \draw [arrow] (pro_pow) |- (out);

    \draw [arrow] (dec_div) -- node[anchor=east] {non} (pro_div);
    \draw [arrow] (dec_div) -- node[anchor=south] {oui} (out_div_err);
    \draw [arrow] (pro_div) |- (out);

    \draw [arrow] (out) -- (stop);
    \draw [arrow] (out_div_err) |- (stop);
    \end{tikzpicture}
}
\end{center}
\subsection*{Programme Python}
\noindent
\begin{minipage}{\linewidth}
\begin{lstlisting}
print("Calculatrice simple (+, -, *, /, **)")
op = input("Votre choix d'opération : ")

if op in ['+', '-', '*', '/', '**']:
    x = float(input("Entrez x : "))
    y = float(input("Entrez y : "))

    if op == '+':
        print(f"Résultat : {x + y}")
    elif op == '-':
        print(f"Résultat : {x - y}")
    elif op == '*':
        print(f"Résultat : {x * y}")
    elif op == '/':
        if y != 0:
            print(f"Résultat : {x / y}")
        else:
            print("Erreur : Division par zéro.")
    elif op == '**':
        print(f"Résultat : {x ** y}")
else:
    print("Erreur : Opération invalide.")
\end{lstlisting}
\end{minipage}

\subsection*{Exemples d\'exécution}
\begin{verbatim}
# Cas 1: Addition (x=10, y=5)
...
Résultat : 10.0 + 5.0 = 15.0
\end{verbatim}
\end{exoBox}

\end{document}